\chapter{Motor neurone disease (MND)}

\section*{Definition and Overview}
Motor neurone disease (MND) is a group of progressive, neurodegenerative disorders characterised by the selective loss of upper and lower motor neurones in the brain, brainstem, and spinal cord. The destruction of these nerve cells prevents the transmission of signals from the brain to the muscles, leading to muscle wasting, weakness, and loss of voluntary movement. Amyotrophic lateral sclerosis (ALS) is the most common form. MND is a terminal condition with no known cure, requiring multidisciplinary supportive care.

\section*{Epidemiology}
MND is a relatively rare but devastating condition. In Australia, the lifetime risk of developing MND is approximately 1 in 300 for men and 1 in 400 for women. It is most commonly diagnosed in individuals aged 50 to 70, although it can occur at any age. Males are slightly more frequently affected than females (ratio of 1.2--1.5:1). About 90--95\% of cases are sporadic (occurring without a family history), while 5--10\% are familial (inherited), often associated with specific genetic mutations.

\section*{Aetiology and Pathophysiology}
The exact cause of sporadic MND remains unknown, but it is believed to result from a complex interaction of genetic, cellular, and environmental factors:

\begin{itemize}
    \item \textbf{Genetic mutations:} Familial MND is linked to mutations in specific genes, such as SOD1, C9orf72, TARDBP, and FUS, which lead to protein misfolding and cellular toxicity.
    \item \textbf{Glutamate excitotoxicity:} Excess accumulation of the neurotransmitter glutamate in the synaptic cleft overstimulates motor neurones, leading to calcium influx and cellular damage.
    \item \textbf{Mitochondrial dysfunction and oxidative stress:} Impaired energy production and accumulation of reactive oxygen species within motor neurones accelerate cell death.
    \item \textbf{Protein aggregation:} The accumulation of abnormal protein aggregates (specifically TDP-43) inside motor neurones disrupts normal cellular transport and function.
    \item \textbf{Upper vs. Lower motor neurones:} Loss of upper motor neurones in the motor cortex leads to muscle spasticity and hyperreflexia, while loss of lower motor neurones in the spinal cord and brainstem causes muscle atrophy, weakness, and fasciculations (twitches).
\end{itemize}

\section*{Clinical Features}
MND symptoms typically start gradually in a localised muscle group (limb or bulbar region) and progress to involve other areas:

\begin{itemize}
    \item \textbf{Limb-onset symptoms:} Weakness in hands or legs, resulting in difficulty gripping objects, writing, buttoning clothes, tripping, or foot drop.
    \item \textbf{Bulbar-onset symptoms:} Difficulty speaking (dysarthria), hoarse voice, and difficulty swallowing (dysphagia), which are common first signs in about 25\% of cases.
    \item \textbf{Muscle wasting and twitches:} Visible muscle atrophy, muscle stiffness (spasticity), and fasciculations (muscle twitching) under the skin.
    \item \textbf{Respiratory symptoms:} Shortness of breath (especially when lying flat), morning headaches, and fatigue due to weakness of the diaphragm and respiratory muscles.
    \item \textbf{Cognitive changes:} Approximately 10--15\% of patients develop frontotemporal lobar degeneration (FTLD), presenting as changes in personality, behaviour, or language.
\end{itemize}

\section*{Differential Diagnosis}
Several other neurological conditions that present with muscle weakness and wasting must be ruled out before diagnosing MND:

\begin{itemize}
    \item \textbf{Cervical spondylotic myelopathy:} Compression of the spinal cord in the neck, causing spasticity in the legs and weakness/sensory changes in the hands.
    \item \textbf{Multifocal motor neuropathy (MMN):} An autoimmune disease causing progressive, asymmetric limb weakness, which can mimic MND but responds to intravenous immunoglobulin (IVIg).
    \item \textbf{Myasthenia gravis:} A neuromuscular junction disorder causing fatiguable muscle weakness, particularly affecting eye movements, speech, and swallowing.
    \item \textbf{Benign fasciculation syndrome:} A common, harmless condition causing muscle twitching without weakness or muscle wasting.
    \item \textbf{Kennedy's disease:} A rare, X-linked genetic disorder causing slowly progressive motor neurone loss, which can be distinguished by genetic testing.
\end{itemize}

\section*{When to Seek Medical Care}
Individuals experiencing progressive muscle weakness, slurred speech, swallowing difficulties, or persistent muscle twitching should consult a GP for urgent referral to a neurologist. MND progresses continuously, and early diagnosis is key to establishing supportive care and accessing clinical trials.

\textbf{Call 000 (or local emergency services) for:}
\begin{itemize}
    \item Sudden, severe difficulty breathing or signs of respiratory failure
    \item Choking on food or saliva that blocks the airway
    \item A fall resulting in severe injury, head trauma, or suspected fracture
    \item Severe chest pain or sudden loss of consciousness
    \item Any acute, life-threatening emergency symptom
\end{itemize}

\section*{Diagnosis and Investigations}
There is no single diagnostic test for MND; diagnosis is based on clinical criteria and ruling out other conditions:

\begin{itemize}
    \item \textbf{Electromyography (EMG) and Nerve Conduction Studies (NCS):} Crucial tests to detect active degeneration of lower motor neurones in multiple body regions and confirm normal sensory nerve function.
    \item \textbf{MRI scan of the brain and spine:} Performed to rule out other structural causes of spinal cord or nerve root compression.
    \item \textbf{Blood tests:} Checking creatine kinase (CK) levels (often mildly elevated in MND) and ruling out thyroid disease, vitamin deficiencies, and heavy metal exposure.
    \item \textbf{Genetic testing:} Offered to patients with a family history of MND to identify known causative mutations.
    \item \textbf{Respiratory function tests:} Regularly measuring vital capacity to monitor respiratory muscle strength.
\end{itemize}

\section*{Management}
The management of MND is palliative and supportive, focusing on symptom control, quality of life, and function:

\begin{itemize}
    \item \textbf{Multidisciplinary care:} Coordinated care involving neurologists, respiratory physicians, occupational therapists, physiotherapists, speech pathologists, dietitians, and palliative care specialists.
    \item \textbf{Disease-modifying drugs:} Riluzole, a glutamate antagonist, is the standard oral medication that can modestly extend survival by 2 to 3 months.
    \item \textbf{Symptom management:} Using baclofen or tizanidine for muscle spasticity, anticholinergic medications (e.g., glycopyrrolate or amitriptyline) for excess saliva (sialorrhoea), and low-dose opioids for breathlessness and anxiety.
    \item \textbf{Respiratory support:} Non-invasive ventilation (NIV)---using a mask overnight---to support weakened breathing muscles, improve sleep, and extend survival.
    \item \textbf{Nutritional support:} Placing a gastrostomy tube (PEG) directly into the stomach when swallowing becomes unsafe or weight loss is severe.
    \item \textbf{Assistive technology:} Communication aids (e.g., eye-gaze technology), wheelchairs, and home modifications to maintain independence.
\end{itemize}

\section*{Complications}
\begin{itemize}
    \item Respiratory failure, the most common cause of death, resulting from diaphragmatic weakness
    \item Aspiration pneumonia caused by food, liquids, or saliva entering the lungs due to bulbar weakness
    \item Malnutrition and severe dehydration due to dysphagia (swallowing difficulties)
    \item Venous thromboembolism (deep vein thrombosis) due to immobility
    \item Pain from muscle spasticity, immobility, and pressure on joints
    \item Depression, anxiety, and severe emotional distress for the patient and their carers
\end{itemize}

\section*{Prognosis and Outcomes}
MND is a progressive and terminal illness with a variable rate of decline. The average survival from the onset of symptoms is 2 to 5 years, although approximately 10\% of patients live for more than 10 years. Bulbar-onset MND typically progresses more rapidly than limb-onset MND. Early introduction of non-invasive ventilation and gastrostomy, alongside coordinated multidisciplinary care, has been shown to improve survival, reduce complications, and significantly enhance the patient's quality of life.

\section*{Prevention and Self-Care}
\begin{itemize}
    \item There are no proven methods to prevent MND, as the precise causes of sporadic cases remain unknown
    \item Engage in gentle, active-assisted range-of-motion exercises under a physiotherapist's guidance to maintain joint flexibility and reduce pain
    \item Adopt a soft, pureed, or high-calorie diet and eat smaller, more frequent meals if swallowing becomes tiring
    \item Use communication boards or text-to-speech apps early in the disease to establish alternative communication methods
    \item Discuss advance care planning with family and doctors early to outline preferences for life-sustaining treatments (like mechanical ventilation)
\end{itemize}

\section*{Sources and Further Reading}
The following organisations provide reliable information on motor neurone disease:

\begin{itemize}
    \item Motor Neurone Disease Australia. \textit{Support services, patient guides, and research updates in Australia.} https://www.mndaustralia.org.au/
    \item Healthdirect Australia. \textit{Overview of MND symptoms, diagnosis, and support networks.} https://www.healthdirect.gov.au/
    \item MND Association (UK). \textit{Comprehensive resources for patients, families, and healthcare professionals.} https://www.mndassociation.org/
    \item Mayo Clinic. \textit{Clinical overview of amyotrophic lateral sclerosis (ALS) and supportive therapies.} https://www.mayoclinic.org/
    \item ALS Association. \textit{Patient education, advocacy, and research updates on ALS.} https://www.als.org/
\end{itemize}
